\documentclass[10pt,a4paper]{amsart}
\usepackage{amsmath,amssymb,amsthm,graphicx}
\usepackage{verbatim}
\newenvironment{reference}{\comment}{\endcomment}
%\newenvironment{reference}{}{}
\newenvironment{slogan}{\comment}{\endcomment}
\newenvironment{history}{\comment}{\endcomment}

% For commutative diagrams we use Xy-pic
\usepackage[all]{xy}

% We use 2cell for 2-commutative diagrams.
\xyoption{2cell}
% Theorem environments.
%
\theoremstyle{plain}
\newtheorem{theorem}[subsection]{Theorem}
\newtheorem{proposition}[subsection]{Proposition}
\newtheorem{lemma}[subsection]{Lemma}

\theoremstyle{definition}
\newtheorem{definition}[subsection]{Definition}
\newtheorem{example}[subsection]{Example}
\newtheorem{exercise}[subsection]{Exercise}
\newtheorem{situation}[subsection]{Situation}

\theoremstyle{remark}
\newtheorem{remark}[subsection]{Remark}
\newtheorem{remarks}[subsection]{Remarks}

\numberwithin{equation}{subsection}

% Macros
%
\def\lim{\mathop{\mathrm{lim}}\nolimits}
\def\colim{\mathop{\mathrm{colim}}\nolimits}
\def\Spec{\mathop{\mathrm{Spec}}}
\def\Hom{\mathop{\mathrm{Hom}}\nolimits}
\def\Ext{\mathop{\mathrm{Ext}}\nolimits}
\def\SheafHom{\mathop{\mathcal{H}\!\mathit{om}}\nolimits}
\def\SheafExt{\mathop{\mathcal{E}\!\mathit{xt}}\nolimits}
\def\Sch{\mathit{Sch}}
\def\Mor{\mathop{\mathrm{Mor}}\nolimits}
\def\Ob{\mathop{\mathrm{Ob}}\nolimits}
\def\Sh{\mathop{\mathit{Sh}}\nolimits}
\def\NL{\mathop{N\!L}\nolimits}
\def\CH{\mathop{\mathrm{CH}}\nolimits}
\def\proetale{{pro\text{-}\acute{e}tale}}
\def\etale{{\acute{e}tale}}
\def\QCoh{\mathit{QCoh}}
\def\Ker{\mathop{\mathrm{Ker}}}
\def\Im{\mathop{\mathrm{Im}}}
\def\Coker{\mathop{\mathrm{Coker}}}
\def\Coim{\mathop{\mathrm{Coim}}}

% Boxtimes
%
\DeclareMathSymbol{\boxtimes}{\mathbin}{AMSa}{"02}

%
% Macros for moduli stacks/spaces
%
\def\QCohstack{\mathcal{QC}\!\mathit{oh}}
\def\Cohstack{\mathcal{C}\!\mathit{oh}}
\def\Spacesstack{\mathcal{S}\!\mathit{paces}}
\def\Quotfunctor{\mathrm{Quot}}
\def\Hilbfunctor{\mathrm{Hilb}}
\def\Curvesstack{\mathcal{C}\!\mathit{urves}}
\def\Polarizedstack{\mathcal{P}\!\mathit{olarized}}
\def\Complexesstack{\mathcal{C}\!\mathit{omplexes}}
% \Pic is the operator that assigns to X its picard group, usage \Pic(X)
% \Picardstack_{X/B} denotes the Picard stack of X over B
% \Picardfunctor_{X/B} denotes the Picard functor of X over B
\def\Pic{\mathop{\mathrm{Pic}}\nolimits}
\def\Picardstack{\mathcal{P}\!\mathit{ic}}
\def\Picardfunctor{\mathrm{Pic}}
\def\Deformationcategory{\mathcal{D}\!\mathit{ef}}
\usepackage{cleveref}
\crefname{originalequation}{original equation}{original equations}
\Crefname{originalequation}{Original equation}{Original equations}
\makeatletter
\newlabel{equation-FM-QCoh}{{0FZC}{source}{}{sourcepointer-1}{}}
\newlabel{equation-FM-QCoh@cref}{{[originalequation][0][]0FZC}{[1][0][]source}{}{}{}}
\makeatother
\crefname{originallemma}{original lemma}{original lemmas}
\Crefname{originallemma}{Original lemma}{Original lemmas}
\makeatletter
\newlabel{lemma-functor-quasi-coherent-from-affine-diagonal}{{0FZH}{source}{}{sourcepointer-2}{}}
\newlabel{lemma-functor-quasi-coherent-from-affine-diagonal@cref}{{[originallemma][0][]0FZH}{[1][0][]source}{}{}{}}
\makeatother
\crefname{originalreference}{original reference}{original references}
\Crefname{originalreference}{Original reference}{Original references}
\makeatletter
\newlabel{coherent-lemma-support-proper-over-base-pushforward}{{08DS}{source}{}{sourcepointer-3}{}}
\newlabel{coherent-lemma-support-proper-over-base-pushforward@cref}{{[originalreference][0][]08DS}{[1][0][]source}{}{}{}}
\makeatother
\crefname{originallemma}{original lemma}{original lemmas}
\Crefname{originallemma}{Original lemma}{Original lemmas}
\makeatletter
\newlabel{lemma-functor-coherent}{{0FZL}{source}{}{sourcepointer-4}{}}
\newlabel{lemma-functor-coherent@cref}{{[originallemma][0][]0FZL}{[1][0][]source}{}{}{}}
\makeatother
\crefname{originallemma}{original lemma}{original lemmas}
\Crefname{originallemma}{Original lemma}{Original lemmas}
\makeatletter
\newlabel{lemma-functor-quasi-coherent-from-separated}{{0FZJ}{source}{}{sourcepointer-5}{}}
\newlabel{lemma-functor-quasi-coherent-from-separated@cref}{{[originallemma][0][]0FZJ}{[1][0][]source}{}{}{}}
\makeatother
\crefname{originalreference}{original reference}{original references}
\Crefname{originalreference}{Original reference}{Original references}
\makeatletter
\newlabel{morphisms-definition-scheme-theoretic-support}{{05JV}{source}{}{sourcepointer-6}{}}
\newlabel{morphisms-definition-scheme-theoretic-support@cref}{{[originalreference][0][]05JV}{[1][0][]source}{}{}{}}
\makeatother
\crefname{originallemma}{original lemma}{original lemmas}
\Crefname{originallemma}{Original lemma}{Original lemmas}
\makeatletter
\newlabel{lemma-characterize-finite}{{0FZM}{source}{}{sourcepointer-7}{}}
\newlabel{lemma-characterize-finite@cref}{{[originallemma][0][]0FZM}{[1][0][]source}{}{}{}}
\makeatother
\title{k-linear exact functors on coherent modules over finite type schemes are FM functors with kernels finite over Y}
\author{The Stacks Project Authors}
\date{stacks-project:functors.tex:local-snapshot-20261012; GNU FDL 1.2 or later}
\begin{document}
\maketitle
\noindent\textit{Editorial note.} All delivered mathematical text is the author's own, copied verbatim from the pinned Stacks Project chapter file functors.tex; no statement or proof step was rewritten, reordered or removed.\par\medskip
\noindent\textit{Editorial note.} Span recomputation by label pairing (LANE\_BRIEF 10.2): in this chapter the author's proof environment follows the statement environment as a standalone block; the spec covers the complete environment pair (statement span, proof span) as recorded in the pool candidate.\par\medskip
\noindent\textit{Editorial note.} Occupation scan: combined claim\_key + statement\_anchor + generic-line scan over all registration files (active and withdrawn) found 0 hits for functors.tex (fu-occupied.json, 2026-10-14); consistent with the duty manager's six-check dispatch.\par\medskip
\ifdefined\hypertarget\hypertarget{sourcepointer-1}{}\fi
\noindent\textit{Source reference.} Original equation 0FZC: sources/stacks-chow/tags.txt, source lines 19947--19947. Cross-chapter standard prerequisite cited by the author inside the counted span; it is used as a pointer and its own proof is not part of this item.\par\medskip
\ifdefined\hypertarget\hypertarget{sourcepointer-2}{}\fi
\noindent\textit{Source reference.} Original lemma 0FZH: sources/stacks-chow/tags.txt, source lines 19952--19952. Cross-chapter standard prerequisite cited by the author inside the counted span; it is used as a pointer and its own proof is not part of this item.\par\medskip
\ifdefined\hypertarget\hypertarget{sourcepointer-3}{}\fi
\noindent\textit{Source reference.} Original reference 08DS: sources/stacks-chow/tags.txt, source lines 10647--10647. Cross-chapter standard prerequisite cited by the author inside the counted span; it is used as a pointer and its own proof is not part of this item.\par\medskip
\ifdefined\hypertarget\hypertarget{sourcepointer-4}{}\fi
\noindent\textit{Source reference.} Original lemma 0FZL: sources/stacks-chow/tags.txt, source lines 19956--19956. Cross-chapter standard prerequisite cited by the author inside the counted span; it is used as a pointer and its own proof is not part of this item.\par\medskip
\ifdefined\hypertarget\hypertarget{sourcepointer-5}{}\fi
\noindent\textit{Source reference.} Original lemma 0FZJ: sources/stacks-chow/tags.txt, source lines 19954--19954. Cross-chapter standard prerequisite cited by the author inside the counted span; it is used as a pointer and its own proof is not part of this item.\par\medskip
\ifdefined\hypertarget\hypertarget{sourcepointer-6}{}\fi
\noindent\textit{Source reference.} Original reference 05JV: sources/stacks-chow/tags.txt, source lines 7179--7179. Cross-chapter standard prerequisite cited by the author inside the counted span; it is used as a pointer and its own proof is not part of this item.\par\medskip
\ifdefined\hypertarget\hypertarget{sourcepointer-7}{}\fi
\noindent\textit{Source reference.} Original lemma 0FZM: sources/stacks-chow/tags.txt, source lines 19957--19957. Cross-chapter standard prerequisite cited by the author inside the counted span; it is used as a pointer and its own proof is not part of this item.\par\medskip
\begin{lemma}
\label{lemma-functor-coherent-over-field}
Let $k$ be a field. Let $X$, $Y$ be finite type schemes over $k$ with
$X$ separated. There is an equivalence of categories between
\begin{enumerate}
\item the category of $k$-linear exact functors
$F : \textit{Coh}(\mathcal{O}_X) \to \textit{Coh}(\mathcal{O}_Y)$, and
\item the category of coherent $\mathcal{O}_{X \times Y}$-modules
$\mathcal{K}$ which are flat over $X$ and have support finite over $Y$
\end{enumerate}
given by sending $\mathcal{K}$ to the restriction of the functor
(\ref{equation-FM-QCoh}) to $\textit{Coh}(\mathcal{O}_X)$.
\end{lemma}

\begin{proof}
Let $\mathcal{K}$ be as in (2). By
Lemma \ref{lemma-functor-quasi-coherent-from-affine-diagonal}
the functor $F$ given by (\ref{equation-FM-QCoh}) is exact and $k$-linear.
Moreover, $F$ sends $\textit{Coh}(\mathcal{O}_X)$ into
$\textit{Coh}(\mathcal{O}_Y)$ for example by
Cohomology of Schemes, Lemma
\ref{coherent-lemma-support-proper-over-base-pushforward}.

\medskip\noindent
Let us construct the quasi-inverse to the construction. Let $F$ be
as in (1). By Lemma \ref{lemma-functor-coherent} we can extend $F$
to a $k$-linear exact functor on the
categories of quasi-coherent modules which commutes with arbitrary direct sums.
By Lemma \ref{lemma-functor-quasi-coherent-from-affine-diagonal}
the extension corresponds to a unique quasi-coherent module
$\mathcal{K}$, flat over $X$, such that
$R^q\text{pr}_{2, *}(\text{pr}_1^*\mathcal{F}
\otimes_{\mathcal{O}_{X \times Y}} \mathcal{K}) = 0$ for $q > 0$
for all quasi-coherent $\mathcal{O}_X$-modules $\mathcal{F}$.
Since $F(\mathcal{O}_X)$ is a coherent $\mathcal{O}_Y$-module, we
conclude from Lemma \ref{lemma-functor-quasi-coherent-from-separated}
that $\mathcal{K}$ is coherent.

\medskip\noindent
For a closed point $x \in X$ denote $\mathcal{O}_x$ the skyscraper sheaf
at $x$ with value the residue field of $x$. We have
$$
F(\mathcal{O}_x) =
\text{pr}_{2, *}(\text{pr}_1^*\mathcal{O}_x \otimes \mathcal{K}) =
(x \times Y \to Y)_*(\mathcal{K}|_{x \times Y})
$$
Since $x \times Y \to Y$ is finite, we see that the pushforward along
this morphism is faithful. Hence if $y \in Y$ is in the image of the
support of $\mathcal{K}|_{x \times Y}$, then $y$ is in the support of
$F(\mathcal{O}_x)$.

\medskip\noindent
Let $Z \subset X \times Y$ be the scheme theoretic support $Z$ of
$\mathcal{K}$, see
Morphisms, Definition \ref{morphisms-definition-scheme-theoretic-support}.
We first prove that $Z \to Y$ is quasi-finite, by proving that its fibres
over closed points are finite. Namely, if the fibre of $Z \to Y$ over a
closed point $y \in Y$ has dimension $> 0$, then we can find infinitely
many pairwise distinct closed points $x_1, x_2, \ldots$ in the image of
$Z_y \to X$. Since we have a surjection
$\mathcal{O}_X \to \bigoplus_{i = 1, \ldots, n} \mathcal{O}_{x_i}$
we obtain a surjection
$$
F(\mathcal{O}_X) \to \bigoplus\nolimits_{i = 1, \ldots, n} F(\mathcal{O}_{x_i})
$$
By what we said above, the point $y$ is in the support of each
of the coherent modules $F(\mathcal{O}_{x_i})$. Since $F(\mathcal{O}_X)$
is a coherent module, this will lead to a contradiction because
the stalk of $F(\mathcal{O}_X)$ at $y$ will be generated by $< n$ elements
if $n$ is large enough. Hence $Z \to Y$ is quasi-finite.
Since $\text{pr}_{2, *}\mathcal{K}$ is coherent and
$R^q\text{pr}_{2, *}\mathcal{K} = 0$ for $q > 0$ we conclude
that $Z \to Y$ is finite by Lemma \ref{lemma-characterize-finite}.
\end{proof}

\end{document}
